\chapter{Holonomic Relation Graphs and Geometry-First Dependencies}
\label{chap:geometry-dependencies}

\section{Pairs are relations, not exclusive partners}

The phrase ``every structure has a pair'' is useful only after the word
\emph{pair} is typed.  The half-axis is paired with binary complement through a
fixed-point relation, with Shannon entropy through an extremal relation, with
the Bloch equator through a representation map, and with destructive
interference through a cancellation theorem.  These are not four copies of one
edge.  They carry different mathematical content.

Accordingly, the v0.7 representation is a graph
\begin{equation}
 \mathcal G=(V,E),
 \label{eq:semantic-relation-graph}
\end{equation}
not a partition of $V$ into disjoint two-element sets.  The requirement is
\begin{equation}
 \deg(v)\ge1
 \qquad\text{for every declared semantic node }v,
 \label{eq:no-isolated-semantic-node}
\end{equation}
while nodes may have several typed neighbours.  This allows one concept to be
related simultaneously to its dual, its representation, its consequence and
its cross-reference.

\section{Relation types}

For the present programme it is sufficient to distinguish at least the
following edge classes:
\begin{description}
 \item[Fixed point.] A map $F$ selects a node by $F(v)=v$.
 \item[Dual.] Two nodes are exchanged by an involution or reciprocal map.
 \item[Implies.] A set of premises logically entails a conclusion.
 \item[Represents.] Two descriptions are coordinates of the same declared
 mathematical object inside a stated representation.
 \item[Holonomic.] A connection/path law assigns a loop or path phase.
 \item[Cross-reference.] A result is imported from another TIR/PhaseNav module
 without claiming that the two objects are identical.
\end{description}

The distinction immediately resolves a recurring ambiguity.  A semantic pair
does not automatically carry a phase.  Holonomy is meaningful only when a
connection and a closed or composable path have been specified.

\section{Normalized holonomy}

At the primitive level, PhaseNav need not store holonomy in radians.  Let an
eligible edge $e$ carry
\begin{equation}
 h(e)\in\mathbb R/\mathbb Z
 \label{eq:normalized-holonomy-edge}
\end{equation}
measured in turns.  A half-turn has $h=1/2$ and a full closure has $h=0$ modulo
one.  If a path
\begin{equation}
 P=(v_0,e_1,v_1,\ldots,e_n,v_n)
\end{equation}
consists of compatible holonomic edges, define
\begin{equation}
 h(P)=\sum_{k=1}^{n}h(e_k)\pmod1.
 \label{eq:path-holonomy-sum}
\end{equation}
A complex $U(1)$ representative may then be introduced as
\begin{equation}
 \mathcal U(P)=\exp(2\pi i\,h(P)).
 \label{eq:u1-turn-representation}
\end{equation}
The representation uses radians; the path composition law does not require
radians as a primitive.

The Berry phase is a standard $U(1)$ connection holonomy on a quantum-state
bundle \citep{berry1984,simon1983}.  The Aharonov--Bohm phase is likewise a
$U(1)$ holonomy of an electromagnetic connection \citep{aharonovbohm1959}.
Their common bundle structure is stronger than a metaphor but weaker than an
identification of physical potentials.  The present monograph uses Berry
holonomy directly; Aharonov--Bohm appears only as a structural cross-reference
to the wider TIR holonomy programme.

\section{The typed half-axis graph}

A minimal graph relevant to this monograph contains the following edges:
\begin{longtable}{@{}p{0.20\textwidth}p{0.22\textwidth}p{0.19\textwidth}p{0.29\textwidth}@{}}
\toprule
Source & Target & Type/status & Meaning\\
\midrule
$\sigma=1/2$ & complement fixedness & fixed point / exact & $\sigma=1-\sigma$.\\
$\sigma=1/2$ & $H_2=\ln2$ & implication / exact & unique entropy maximum.\\
$\sigma=1/2$ & Bloch equator & representation / standard & $n_z=2\sigma-1=0$.\\
Bloch equator & Berry $-1$ & holonomic / standard & equatorial loop, $h=1/2$.\\
$\sigma=1/2$ + half-turn & cancellation zero & implication / exact & equal-gain detector.\\
$q\in\mathbb R/\mathbb Z$ & projective cycle & representation / exact & normalized recurrence.\\
projective cycle & spinor return & holonomic / standard & double cover, sheet flip.\\
centered zeta axis & reciprocal axis & dual / exact & ambient axis set invariance.\\
$12$ projective cycles & $\ln2/12$ per turn & implication / model & information-cycle assignment.\\
$\ln2/12$ + $C=2\pi$ & $\kappa$ & implication / model & conditional angular density.\\
canonical zero state & native closure & implication / open & SOH-C004.\\
\bottomrule
\end{longtable}

\begin{figure}[htbp]
 \centering
 \includegraphics[width=0.88\textwidth]{typed_identity_relation_graph.pdf}
 \caption{A compact view of the typed relation graph.  The dashed logical
 boundary is semantic rather than graphical: an OPEN edge remains OPEN even if
 it connects nodes already embedded in the same figure.}
 \label{fig:typed-relation-graph}
\end{figure}

The graph makes one logical fact visually unavoidable: proximity and
connectivity are not proof.  A path can exist through a semantic graph while an
edge on that path remains a postulate.

\section{Geometry-first dependency direction}

The companion CIEL/PhaseNav branch introduces a low-level package
\begin{center}
 \texttt{ciel\_geometry.phasenav.dependencies}
\end{center}
whose dependency direction is intentionally one-way:
\begin{equation}
\boxed{
\begin{array}{c}
\text{identity axes and centered coordinates}\\
\text{normalized cycles and winding}\\
\text{phase crystal }(\mathbb R/\mathbb Z)^N\\
\text{typed relations with optional holonomy}
\end{array}
\longrightarrow
\text{tangent/log-map vectors}
\longrightarrow
\text{higher PhaseNav operators}.
}
\label{eq:dependency-direction}
\end{equation}
The essential design constraint is that vectorization consumes geometry rather
than generating geometry from opaque identifiers.

Earlier experimental PhaseNav registries can deterministically map a stable
identifier into a high-dimensional vector.  Such a map is useful for a
repeatable computational embedding, but its geometry is not evidence that two
concepts are semantically related.  Version 0.7 therefore separates three
layers:
\begin{enumerate}
 \item declared semantic provenance;
 \item generated candidate geometry;
 \item proof status of any inference that uses the relation.
\end{enumerate}
This separation mirrors the distinction between a numerical embedding and a
mathematical theorem.

\section{The phase crystal}

Let the PhaseNav crystal be
\begin{equation}
 \mathbb T^N=(\mathbb R/\mathbb Z)^N,
 \label{eq:phase-crystal}
\end{equation}
with $N=36$ for the current native geometry.  For two states
$q=(q_1,\ldots,q_N)$ and $r=(r_1,\ldots,r_N)$, define the componentwise shortest
oriented tangent
\begin{equation}
 \Delta_k(q,r)
 =\left(r_k-q_k+\frac12\right)\bmod1-\frac12.
 \label{eq:phase-crystal-logmap}
\end{equation}
Then
\begin{equation}
 \mathbf v(q\to r)
 =(\Delta_1,\ldots,\Delta_N)\in[-1/2,1/2)^N
 \label{eq:phase-crystal-vector}
\end{equation}
is a natural local vectorization operator.  It is analogous to a logarithmic
map in the sense that it converts a pair of points into a tangent displacement,
with the usual caveat that at exactly antipodal coordinates the shortest branch
requires a convention.

This construction generates negative, zero and positive vector components
without adding sign labels after the fact:
\begin{equation}
 \Delta_k<0\;\text{and}\;\Delta_k>0
\end{equation}
are opposite orientations, whereas $\Delta_k=0$ means that the two phase
coordinates coincide.  For a centered binary coordinate, the same principle
reduces to \cref{eq:identity-sign}.

\section{From phase geometry to vector operators}

Once a tangent vector is available, familiar operators become downstream
objects.  For example, a local inner product may be defined by
\begin{equation}
 \langle v,w\rangle_g=v^Tg(q)w,
 \label{eq:local-inner-product}
\end{equation}
where $g$ is a declared metric.  Projection, orthogonalization, curvature
estimation and constrained updates can then be defined on the tangent space.
The important point is architectural:
\begin{equation}
 \text{metric/connection}
 \longrightarrow
 \text{tangent relation}
 \longrightarrow
 \text{vector operator},
\end{equation}
not the reverse.

The same philosophy applies to semantic mass, density or viscosity fields.
Those quantities may depend on geometry, but the fact that two stable IDs happen
to be close under one deterministic embedding cannot establish a semantic law.
An explicit provenance edge is still required.

\section{Holonomy eligibility}

The following rule is enforced conceptually and in the companion runtime:
\begin{criterion}[Holonomy eligibility]
\label{crit:holonomy-eligibility}
A semantic edge may carry canonical holonomy only if
\begin{enumerate}[label=(H\arabic*)]
 \item a connection or equivalent phase-transport law is specified;
 \item the edge belongs to a closed loop or a composable path with a defined
 endpoint convention;
 \item the resulting phase is invariant under the permitted gauge changes, or
 the gauge dependence is explicitly part of the object;
 \item its claim status is independent of a purely generated semantic
 embedding.
\end{enumerate}
\end{criterion}

Hence the pair ``energy--matter'' does not automatically receive a $\pi$ phase,
and neither does ``time--space''.  In the wider TIR interpretation these may be
paired sectors, but holonomy requires additional geometric data.

\section{Arithmetic and phase geometry as paired sectors}

The current TIR crosswalk proposes the model-level split
\begin{equation}
 \text{information arithmetic}\longrightarrow\text{energy and matter},
 \label{eq:tir-arithmetic-sector}
\end{equation}
paired with
\begin{equation}
 \text{phase geometry}\longrightarrow\text{time and space}.
 \label{eq:tir-phase-sector}
\end{equation}
This monograph does not use either arrow as an established law of physics.
They are recorded because they define the intended dependency architecture of
the broader programme.  In particular, a proposed derivation of physical time
$t$ from a phase period $\tau$ is deferred.  No theorem in the present work
requires it.

The useful part for Secret of a Half is more modest.  The native theta bridge
already separates $\delta=\Re s-1/2$ from phase transport $t$.  The
geometry-first dependency package provides a consistent way to represent that
separation without conflating transverse identity with tangent evolution.

\section{TIR $W_{ij}$ and scope boundary}

The wider Theory of Informational Relations contains a relation object commonly
denoted $W_{ij}$.  In the current work it is used only as a cross-reference to a
broader holonomic relation/information architecture.  No definition, equation
of motion or physical identification of $W_{ij}$ is imported into the proof
chain here.  If a future version uses $W_{ij}$ as an executable edge weight, it
must satisfy \cref{crit:holonomy-eligibility} and enter the claim ledger with an
explicit status.

\section{Native PhaseNav declaration}

The v0.7 semantic source
\begin{center}
 \texttt{construction/phasenav/secret\_of\_half\_identity\_holonomy\_v0\_7.pnv}
\end{center}
records the nodes, normalized holonomy turns, cross-factorizations, TIR
cross-references and firewall status.  The declaration stores a half-turn as
$1/2$ turn rather than as the primitive number $\pi$; the radian conversion is a
representation step.

This file does not supersede the theta-bridge, Weil, Hermite, reciprocal-tail or
DHSE constructions.  It is an integration layer.  Existing numerical receipts
remain untouched.

\section{Claim boundary}

The phase-crystal tangent formula, normalized path composition and semantic
graph rules are exact definitions.  Berry and spinor holonomies use standard
geometry and representation theory.  The TIR sector assignments, the
information-cycle interpretation of twelve, and any physical identification of
semantic edges remain model-level.

\status{Version 0.7 introduces a geometry-first dependency discipline and a
typed semantic relation graph.  Generated PhaseNav geometry is diagnostic until
independent provenance promotes the relevant relation.}
